
\section{Descent of theta periods through their actual matching quotients}
\label{sec:sr-theta-quotient-descent}

The square-integrability statement just proved supplies the positive
Hermitian form needed for descent. Merely knowing a one-dimensional
space of local branching maps would not supply this step.

\begin{srlemma}[A unitary image kills the radical of a local theta lift]
\label{lem:sr-theta-unitary-head}
Let \(Q\) be a finite-length admissible smooth representation of a
nonarchimedean reductive group, or an admissible finite-length
Harish--Chandra module for a real reductive group. Suppose its maximal
semisimple quotient is an irreducible \(\pi\), with multiplicity one.
A map from \(Q\) into a unitary representation factors through
\(Q\longrightarrow\pi\), provided in the real case it intertwines
the compact group and the infinitesimal actions on smooth vectors.
\end{srlemma}

\begin{proof}
Give its image \(M\) the inherited positive Hermitian form.
For an invariant submodule \(N\subset M\), take closure in the Hilbert
space. In the nonarchimedean case, a vector fixed by a compact open
group \(K\) projects into the closure of \(N^K\); the latter space is
finite dimensional by admissibility and hence closed. Thus its
orthogonal projection belongs to \(N\). In the real case perform the
same argument on the finitely many compact types of a given vector.
Every such isotypic component of \(N\) is finite dimensional by
admissibility and closed, so again the projection belongs to \(N\).
It follows that
\(M=N\oplus(N^\perp\cap M)\) algebraically.
The complement is invariant under the group in the first case.
In the real case it is invariant under the compact group and under
the Lie algebra, since the infinitesimal unitary action is skew
adjoint on these smooth vectors.
Induction on length makes \(M\) semisimple. Every semisimple quotient
of \(Q\) factors through its maximal semisimple quotient, proving
the claim. This argument allows the map to be zero.
\end{proof}

We recall exactly the classical local theorem used with this lemma.
For a characteristic-zero local unitary dual pair and an irreducible
admissible representation of its first member, the big theta lift is
admissible of finite length and its maximal semisimple quotient is
irreducible or zero. At real places Howe's Theorem~1A (p.~536) gives the finitely
generated admissible quasisimple big lift and its unique irreducible
quotient; its polynomial algebraic version is Theorem~2.1
(pp.~538--539). At nonsplit finite places it is Howe duality in
arbitrary residual characteristic.
At a split finite place the pair is the corresponding type II
general linear pair, with the same finite-length and unique-head
statement. We use respectively Howe,
\emph{Transcending classical invariant theory};
Gan--Takeda, \emph{A proof of the Howe duality conjecture}, including
the finite-length assertion recalled before their main theorem;
and M\'inguez, \emph{Correspondance de Howe explicite: paires duales
de type II}, Corollary~6.3 and the big-lift finite-length statement.
\footnote{Primary texts:
\url{https://www.math.utah.edu/~ptrapa/math-library/howe/howe-transcending-1989.pdf},
\url{https://arxiv.org/abs/1407.1995},
\url{https://arxiv.org/abs/0709.4384}.}

Their application here is to the trivial representation of the
\emph{actual first group}, after its specified splitting.
Thus the real polynomial realization has character parameters
\(\det^c\) and \(\det^{-c}\) in the canonical first Fock actions,
with \(c=(m_\alpha-m_\beta)/2\), exactly as calculated above.
The target twist is the already checked \(\det^c\) on the primed
canonical second action. The finite pairs are the actual
\(U(X_v),U(W_v)\), or their displayed split general linear models.
All local fields have characteristic zero, including places above
\(2\). The genuine covers, trace character, and splitting twists
are those already used to construct the matching maps.
The nonzero maps constructed there identify the small local theta
lift with the specified \(\Pi_v\), or \(\Pi_\infty\); no new
local quotient is chosen.

\begin{srproposition}[Actual period factorization]
\label{prop:sr-theta-actual-period-factor}
On polynomial infinite vectors, the large theta period factors through
\[
       \Pi=\Pi_\infty\otimes
                   \bigotimes\nolimits'_v\Pi_v .
 \tag{D1}\label{eq:sr-theta-common-global-module}
\]
The small theta period, on its actual full-rank cutoff, factors through
the same space, using its already chosen primed matching maps.
In particular a polynomial matching relation with zero on one side
has zero theta period on that side. These are exactly the two
vanishing assertions of the source condition \(\mathrm{Factor}\).
\end{srproposition}

\begin{proof}
Fix every local test factor except one.
The resulting theta-period map has values in
\(L^2(U(W)(F)\backslash U(W)(\mathbb A_F))\) by the preceding section.
Right translation on this space is unitary and agrees with the
actual second oscillator action. The theta period is invariant
under the first unitary action: it is just a change of variable
in its convergent first quotient integral. Unimodularity gives
invariance of its Haar measure. The dominated Schwartz estimates
justify differentiation of this identity.

At a finite place other than the selected cutoff place, the local
map therefore factors first through the big theta lift of the
trivial first representation. Its unitary image factors through
the small lift by the preceding lemma, hence through the actual
chosen matching map \(p_v\), or \(p'_v\).
At infinity the same argument applies to the polynomial
Harish--Chandra module. All required differentiations are valid in
\(L^2\): retaining the Schwartz seminorm in the square-integrability
proof gives
\(\|I(f\otimes\phi)\|_2\le C p(f)\) for each fixed finite test space.
The smooth oscillator action and its differentiated polynomial
action are consequently carried into the genuine smooth and
infinitesimal right actions in \(L^2\). The matching quotient is the
already constructed real theta quotient, so its kernel is killed.

There is a separate useful point at the small selected split place
\(v_1\). Here varying a factor must keep it in the actual cutoff.
In its opposite linear model this is the Schwartz space supported
on the open full-row-rank frame locus.
Its first-group coinvariant quotient is already the actual
irreducible \(\Pi_{v_1}\), as proved in the good split coinvariant
calculation. Explicitly, the full-rank frame bundle has free
\(\mathrm{GL}_{a'}\) fibers over a Grassmannian, and fiber integration
with the specified character identifies its coinvariants with the
common normalized parabolic induction. That induction is irreducible.
Its map onto the whole Schwartz coinvariant quotient is nonzero
and hence an isomorphism. Thus its kernel is precisely its own
first augmentation space. Invariance of the convergent small period
kills this kernel directly; the map at \(v_1\) also descends to the
actual \(\Pi_{v_1}\). This argument never evaluates the small
period on vectors outside its convergence cutoff.

These local factorizations imply the tensor assertion by the
universal property of the algebraic tensor product. For a fixed
finite set of nonvacuum places, the theta period is multilinear
in its infinite factor and these finitely many finite factors.
It kills the kernel of each quotient map, so it factors through
their tensor product. Enlarging this set inserts the same nonzero
matched vacuum and gives the restricted product (D1).

For clarity, a displayed finite sum representing a small cutoff
vector may initially have factors outside the cutoff at \(v_1\).
Use a linear projection of their finite dimensional span onto the
subspace satisfying the cutoff. Every infinite slice of the
displayed sum belongs to that subspace, so applying the projection
to the sum leaves the actual Schwartz vector unchanged.
The finite tensor embedding into functions is injective:
evaluate successively against a linearly independent finite basis
of functions in each variable. Hence two equal finite tensor
presentations have the same image under the matching tensor
quotient. We can therefore apply the preceding factorization to
the projected presentation without altering any polynomial
matching relation. This proves the asserted zero-relation
vanishing for both sides.
\end{proof}

\begin{srlemma}[Nonzero periods are injective on the common module]
\label{lem:sr-theta-nonzero-period-injective}
The common module (D1) is algebraically irreducible for its finite
local group actions and infinite Lie and compact actions.
Consequently each of the descended theta-period maps is either
zero or injective.
\end{srlemma}

\begin{proof}
Each local factor is irreducible and has scalar endomorphism ring.
At finite places this follows from admissible Schur's lemma; at
infinity it is the scalar endomorphism assertion already proved
for the countable polynomial quotient.
The algebraic density theorem then says that finitely many
linearly independent vectors in one factor can be sent to any
prescribed vectors by an element of its action algebra.
It follows that a finite tensor product of these factors is
irreducible: apply these operators to isolate one nonzero pure
term of a finite tensor sum, then generate every pure tensor.
Every vector in the restricted product belongs to such a finite
tensor product. Acting at further places generates each local
factor from its nonzero vacuum. Thus a nonzero submodule contains
every restricted pure tensor. The kernel of a descended period
is stable under these actions, proving the last assertion.
\end{proof}

\begin{srlemma}[Local cone augmentation, including the cutoff place]
\label{lem:sr-theta-cone-augmentation}
Let \(D_v\) be either actual local cone operator, with the exact
exponent \(z_A=n-d/2\) or \(z_B=1+d/2-n\).
For every lower first-group invariant linear functional
\(\ell\) on the finite lower Schwartz space,
\[
    \ell(D_v\omega_v(g,1)\phi)=\ell(D_v\phi)
               \qquad(g\in U(X)(F_v)).
 \tag{D2}\label{eq:sr-theta-cone-augmentation}
\]
In particular \(p_{1,v}D_v\) kills the upper first augmentation
space. This assertion holds on the whole local Schwartz space;
it therefore also compares a cutoff representative with the
ordinary vacuum when their actual upper matching classes agree
at a good split place.
\end{srlemma}

\begin{proof}
Let \(Q_x=N_X\rtimes U(X_1)\) be the stabilizer of the nonzero
isotropic vector \(x\). Set
\[
       F_\phi(g)=
            \ell(e_\epsilon\rho(u_\epsilon(1))
                                      \omega_v(g,1)\phi).
\]
The evaluation invariances, the same corner commutators as in (T9),
and invariance of \(\ell\) show that \(F_\phi\) is left invariant
under \(Q_x(F_v)\). The stabilizer \(Q_x\) is unimodular: its
unipotent part is unimodular and the determinant of the
\(U(X_1)\) action on that part has local absolute value \(1\).
Hence \(Q_x\backslash U(X)\) has invariant quotient measure.

In Iwasawa coordinates for this quotient its integral is, up to
a fixed nonzero Haar scalar,
\[
 \int_{K_{X,v}}\int_{E_v^\times}
       F_\phi(m_X(t,1)k)|t|_{E_v}^{1-n}\,d^\times t\,dk .
\]
Insert the exact two rows of (T9).
For \(A\), substitute \(l=t^{-1}\); for \(B\), put \(l=t\).
The resulting expression is exactly \(\ell(D_v\phi)\), up to
that same Haar scalar. All these integrals are absolutely
convergent. This follows from the finite-image and finite-level
argument for the finite cone operator: \(\ell\) is bounded on
that finite dimensional image, the corner parameter is bounded
at the finite place, and its tail at zero has exponent \(z>0\).
The compact average has finite orbit on a smooth vector.
Right invariance of the quotient measure now proves (D2).

Apply (D2) with every algebraic linear functional on the lower
matching quotient composed with its quotient map. Their
invariance follows because that quotient kills lower first
augmentation. These functionals separate points, so
\(p_{1,v}D_v\) kills the upper augmentation.
At a good split place the whole upper matching kernel equals
that augmentation space. This proves the final assertion as well.
\end{proof}

\begin{srproposition}[Descent of the actual constant-term operators]
\label{prop:sr-theta-constant-term-descent}
Consider either codimension-two step, and suppose the lower
large and small periods are nonzero. This is the nonvanishing
available at the earlier stage of the signature induction.
Each actual zero evaluation or cone evaluation appearing in the
four constant-term formulas descends through the upper matching
quotient to the lower matching quotient. Its descended map
intertwines the actual line parabolic of \(U(W)\), with the
character and exponent of that constant-term component.

These are actual descended operators. Their kernel property is
a consequence of the global constant-term formula and the
injective lower period; it is not assumed as an additional
matching condition.
\end{srproposition}

\begin{proof}
Let \(D\) be one of the actual adelic evaluation operators, including
its actual compact cone average when appropriate.
Its finite and infinite factorization was proved above.
If an upper polynomial vector has zero common matching class,
its upper theta period is zero by the period factorization
proposition. Hence both of its distinct Levi constant-term
components are zero. The corresponding formula (T10), (T11),
or its zero-term counterpart gives
\[
                      I_1(D\Phi)=0.
\]
The lower period is injective on its common module by the preceding
lemma and lower nonvanishing. Therefore the lower matching class
of \(D\Phi\) is zero. This proves the adelic kernel property.

To obtain the local maps, fix pure factors at all other places
whose images after \(D\) are nonzero in the corresponding lower
quotients. Such choices are simultaneous: outside the finite bad
set use the exact nonzero vacuum factors; at each bad split place
use the already calculated standard vector with nonzero lower
image. At infinity ordinary zero evaluation is onto on polynomial
Gaussians. For a cone use the proved compact cone bump: the
nonzero lower period gives a nonzero invariant continuous
functional, whose composition with the actual infinite cone
operator is nonzero. Polynomial Gaussian density and the
continuous seminorm estimates give a polynomial choice with
nonzero lower image.

For a small cone the selected cutoff place needs one extra check.
Its ordinary vacuum has nonzero lower image, by its nonzero
unramified cone factor. The full-rank coinvariant calculation
gives a cutoff vector with the same actual upper matching class
as this vacuum. By (D2) their lower images under \(D_v\) agree.
This supplies a nonzero cutoff factor as well. Thus every
required outside pure tensor can indeed be chosen inside the
allowed global domain.

If a local input has zero upper matching class, pair it with
these fixed outside inputs. The adelic kernel property makes
its output tensor zero. The outside output tensor is nonzero,
so the local lower output class is zero. Hence the actual local
operator factors through the actual upper quotient at every
place. The same argument at infinity uses the algebraic
polynomial quotient and gives its descended infinite map.

Finally right translating a global theta period by the line
parabolic translates its constant term by that same parabolic.
Its unipotent radical acts trivially on the constant-term
component. On its Levi the component has exactly the character
and norm power calculated above, and the lower unitary group
acts by right translation in the lower period. Apply the
injectivity of that lower period and separate the nonzero
outside tensor factors as just done. This proves the stated
local parabolic equivariance for the descended operators.
\end{proof}

The last proposition is an induction step with a fully stated
earlier-stage hypothesis. The square integrability and the
period factorization propositions preceding it hold for every
actual triple, independently of that induction. They remove
the quotient-kernel obstruction from the rank-step comparison;
they do not alone prove that two embeddings of the common
global module into automorphic functions are proportional.
That is the remaining global comparison.
